\documentclass[11pt]{article}
\usepackage[a4paper,margin=24mm]{geometry}
\usepackage[T1]{fontenc}
\usepackage{amsmath,amssymb,amsthm,mathtools,xurl}
\usepackage[hidelinks,hypertexnames=false]{hyperref}
\allowdisplaybreaks\setlength{\emergencystretch}{3em}
\newcommand{\C}{\mathbb C}\newcommand{\R}{\mathbb R}
\newcommand{\tr}{\operatorname{tr}}\newcommand{\norm}[1]{\left\lVert#1\right\rVert}
\newtheorem{theorem}{Theorem}
\title{The original degree-volume contraction receives a uniform heat bound}
\author{}\date{20 September 2026 --- additional SZ-20260920-045 result}
\begin{document}\maketitle

This finite calculation returns the metric heat estimate EM29--30
\cite{EM} to the original degree filtration and the original Toda
notation of \cite{TV}. Its parameter \(v_{N+1}\), not a new metric,
is the exact adjacent heat budget. All required finite identities are
proved below. The full main045 proof and frozen043 remain unchanged.
Connes's Gaussian heat test \cite{Connes} and the spectral-action
context of Connes--van Suijlekom \cite{CV} motivate the observations;
none of their infinite-dimensional theorems is used.

\section{The fixed original cyclic space and its actual degree metrics}

Retain the nonempty original packet of actual zeros, closed under
conjugation and \(\rho\mapsto1-\rho\) with complete orders, its divisor
\(h\), \(g=2\xi\), and \(\mathcal MF_h=g/h\).
The cyclic quotient, physical image and measure are
\begin{equation}\tag{VH1}\label{VH1}
\begin{gathered}
C=\C[S]/(\chi),\quad q=\deg\chi,\quad c=k/2,\quad
M=(M_S-cI)/i,\quad \upsilon=j_h(g/h),\\
\eta[P]=\upsilon^{\otimes k}P(v_1+\cdots+v_k),\quad
\eta[1]=\upsilon^{\otimes k},\\
w_h(u)=|(g/h)(1/2+iu)|^2/(2\pi),\quad m_k=w_h^{*k},\quad
d\mu_\theta(u)=e^{\theta u}m_k(u)\,du,
\quad\theta\in\R,\quad|\theta|<\pi/2.
\end{gathered}
\end{equation}
The parameter \(\theta\) is held fixed during each degree comparison.
The original exponential source bound gives finite moments of every
order throughout this tilt interval, with differentiation locally
dominated there. The mass at zero tilt is \((\int w_h)^k\), not one.
The unit \(\upsilon\) is invertible because its constant jet at each
original zero is the nonzero leading coefficient after exact division
by its full order. In a truncated local ring its inverse is a finite
geometric series. The full annihilator \(\chi\) retains every root
and its original collision exponent; hence \(\eta\) is injective
onto the same cyclic physical image. No root is changed in this note.

Let \(U_j(u;\theta)\) be the real monic orthogonal polynomials of
\(\mu_\theta\), with full squared norms \(\omega_j(\theta)>0\).
Put \(p_j(S)=i^jU_j((S-c)/i)\), and \(b_j=[p_j]_\chi\), in the
fixed remainder basis \(1,S,\ldots,S^{q-1}\). For \(N\ge q-1\),
\begin{equation}\tag{VH2}\label{VH2}
E_N=\left[\frac{i^{-j}b_j}{\sqrt{\omega_j}}\right]_{j=0}^N,
\quad G_N=(E_NE_N^*)^{-1},\quad R_N=E_N^*G_N,\quad
V_N=\det G_N.
\end{equation}
The first \(q\) monic columns are independent. The identities
\(E_NR_N=I\), \(R_N^*R_N=G_N\) follow by multiplication.
Every other lift is \(R_Nu+z\), \(E_Nz=0\), and has squared norm
\(u^*G_Nu+\norm z^2\), since its cross terms vanish by
\(R_N^*z=G_NE_Nz=0\). Thus \(G_N\) is precisely the attained
original metric. Its physical source representative is
\[
\left[\sum_{j=0}^N(R_Nu)_j\,
\frac{U_j((S-c)/i)}{\sqrt{\omega_j}}\right]
(\mathcal D_1+\cdots+\mathcal D_k)F_h^{\otimes k},
\]
where the square bracket denotes the displayed polynomial in \(S\),
evaluated at the original sum \(\mathcal D_j=-x_j\partial_{x_j}\).
Its class is \(\eta u\); its tilted norm is \(u^*G_Nu\).
At every degree the physical unit remains \(\eta[1]\).

Append the actual next column, including its complex phase:
\begin{equation}\tag{VH3}\label{VH3}
a=\frac{i^{-(N+1)}b_{N+1}}{\sqrt{\omega_{N+1}}},\quad
G_{N+1}^{-1}=G_N^{-1}+aa^*,\quad
\alpha=a^*G_Na=\frac{b_{N+1}^*G_Nb_{N+1}}{\omega_{N+1}},\quad
\kappa=1+\alpha=\frac{V_N}{V_{N+1}}.
\end{equation}
The first update is the column sum in VH2. To prove the last equality,
use \(\det(I+xy^*)=1+y^*x\): in a basis beginning with \(x\ne0\)
the matrix is triangular except its first row, and its determinant is
the stated first diagonal entry. The case \(x=0\) is immediate.
Apply it to \(G_N^{1/2}aa^*G_N^{1/2}\) and take determinants.

\section{The complete rank-one path and its physical coordinate map}

Abbreviate \(G=G_N\). For \(0\le\lambda\le1\), retain the exact path
\begin{equation}\tag{VH4}\label{VH4}
E_\lambda=[E_N,\sqrt\lambda\,a],\quad
G_\lambda=(G^{-1}+\lambda aa^*)^{-1}
=G-\frac{\lambda}{1+\lambda\alpha}Gaa^*G,\quad
R_\lambda=E_\lambda^*G_\lambda.
\end{equation}
Multiplying the right-hand inverse formula by \(G^{-1}+\lambda aa^*\)
proves it, including \(\lambda=0\). The same fibre-minimum proof as
VH2 applies to \(E_\lambda\), which is onto. This is a specified
interpolation between the actual degree metrics, not a claim that a
fractional degree exists. Its endpoint lifts are the actual ones.
Moreover
\begin{equation}\tag{VH5}\label{VH5}
\det G_\lambda=\frac{V_N}{1+\lambda\alpha},\quad
G_\lambda'=-G_\lambda aa^*G_\lambda,\quad
a^*G_\lambda a=\frac{\alpha}{1+\lambda\alpha}.
\end{equation}
Differentiate \(G_\lambda^{-1}G_\lambda=I\) for the middle formula;
the others follow from VH3--4. These identities keep the actual
volume, including any original mass factor.

If \(\alpha=0\), positivity of \(G\) gives \(a=0\), and the entire
path is constant. If \(\alpha>0\), define the exact \(G\)-orthogonal
rank-one projection
\(P=aa^*G/\alpha\). It obeys \(P^2=P\), \(P^{\dagger_G}=P\).
Then
\begin{equation}\tag{VH6}\label{VH6}
\begin{gathered}
C_\lambda=G^{-1}G_\lambda=I-P+(1+\lambda\alpha)^{-1}P,\\
S_\lambda=C_\lambda^{1/2}=I-P+(1+\lambda\alpha)^{-1/2}P,\quad
S_\lambda^*GS_\lambda=G_\lambda,\quad
\log C_\lambda=-\log(1+\lambda\alpha)P.
\end{gathered}
\end{equation}
The formulas follow by restricting to the orthogonal range and kernel
of \(P\). Thus \(S_\lambda:C_{G_\lambda}\to C_G\) is an isometry.
It carries all the original physical data by
\begin{equation}\tag{VH7}\label{VH7}
M_\lambda=S_\lambda MS_\lambda^{-1},\quad
\eta_\lambda=\eta S_\lambda^{-1},\quad
\mathbf1_\lambda=S_\lambda[1],\quad
\eta_\lambda\mathbf1_\lambda=\eta[1],\quad
\eta_\lambda M_\lambda S_\lambda=\eta M.
\end{equation}
No transformed marked vector is identified with \([1]\) without this
map. Writing \(r_\lambda=(1+\lambda\alpha)^{-1/2}\), the complete
complex block formula is
\begin{equation}\tag{VH8}\label{VH8}
\begin{split}
M_\lambda={}&PMP+r_\lambda PM(I-P)
+r_\lambda^{-1}(I-P)MP+(I-P)M(I-P),\\
D_\lambda:=M_\lambda-M={}&(r_\lambda-1)PM(I-P)
+(r_\lambda^{-1}-1)(I-P)MP,\\
M_\lambda^{\dagger_G}M_\lambda-M^{\dagger_G}M
={}&M^{\dagger_G}D_\lambda+D_\lambda^{\dagger_G}M
+D_\lambda^{\dagger_G}D_\lambda.
\end{split}
\end{equation}
These follow by multiplication and preserve both complex mixed blocks
and the quadratic term. In particular \([P,M]=0\) makes the heat
observation constant, even if the volume decreases strictly.

\section{Exact heat difference and the adjacent Toda-volume budget}

For a positive metric \(H\) on the same cyclic space write
\(X^{\dagger_H}=H^{-1}X^*H\). Define
\begin{equation}\tag{VH9}\label{VH9}
\Phi_s=\tr e^{-sM^2},\quad
\Psi_s(\lambda)=\tr e^{-sM^{\dagger_{G_\lambda}}M},\quad
\Delta_s(\lambda)=\Re\Phi_s-\Psi_s(\lambda),\quad s>0.
\end{equation}
The first observation is exactly fixed throughout; its full value is
\(\sum_\zeta e_\zeta e^{s(\zeta-c)^2}\) for the roots of \(\chi\)
with all their multiplicities. A triangular basis proves the trace
formula termwise from the convergent exponential series. On a Jordan
block the nilpotent derivatives remain in the full operator; only
their traces vanish. No numerical location of roots is needed.

Let \(H_\lambda=M^{\dagger_{G_\lambda}}M\) in the original coordinates,
and set
\begin{equation}\tag{VH10}\label{VH10}
B_\lambda=M e^{-sH_\lambda}M^{\dagger_{G_\lambda}}
-H_\lambda e^{-sH_\lambda},\qquad \tr B_\lambda=0.
\end{equation}
Then the exact original-coordinate heat derivative and its integral are
\begin{equation}\tag{VH11}\label{VH11}
\begin{split}
\Psi_s'(\lambda)&=s\,a^*G_\lambda B_\lambda a
=\frac{s}{\omega_{N+1}}b_{N+1}^*G_\lambda B_\lambda b_{N+1},\\
\Psi_s(N+1)-\Psi_s(N)
&=\frac{s}{\omega_{N+1}}\int_0^1
b_{N+1}^*G_\lambda B_\lambda b_{N+1}\,d\lambda,\\
\Delta_s(N+1)-\Delta_s(N)&=-[\Psi_s(N+1)-\Psi_s(N)].
\end{split}
\end{equation}
To prove the derivative without a commutativity assumption, VH5 gives
\((M^{\dagger_{G_\lambda}})'=[aa^*G_\lambda,M^{\dagger_{G_\lambda}}]\).
Consequently
\(H_\lambda'=aa^*G_\lambda H_\lambda-
M^{\dagger_{G_\lambda}}aa^*G_\lambda M\).
Differentiating \(\tr H_\lambda^n\) produces \(n\) terms equal by
cyclicity. Uniform convergence on the compact interval justifies
summing the differentiated exponential series, and gives
\((\tr e^{-sH_\lambda})'=-s\tr(e^{-sH_\lambda}H_\lambda')\).
Substitute the preceding formula and cycle factors to obtain VH11.
The zero trace in VH10 follows from the same cyclicity.

Both terms in VH10 are positive \(G_\lambda\)-selfadjoint and at most
\(I/(es)\). Indeed a singular-value decomposition in the exact
isometric coordinates \(G_\lambda^{1/2}\) gives their eigenvalues
\(x e^{-sx}\), \(x\ge0\), and their zeros; the maximum of this scalar
function is \(1/(es)\), by differentiation. Thus
\(-I/(es)\preceq_{G_\lambda}B_\lambda\preceq_{G_\lambda}I/(es)\),
where these inequalities mean inequalities of the stated quadratic forms.
VH5 and VH11 prove
\begin{equation}\tag{VH12}\label{VH12}
\boxed{\quad
|\Psi_s(N+1)-\Psi_s(N)|
\le\frac1e\log(1+\alpha)
=\frac1e\log\frac{V_N}{V_{N+1}}
=\frac{v_{N+1}}e,\qquad s>0.\quad}
\end{equation}
The sign of the actual integrand is not prescribed. The volume
contraction is a bound on its absolute integral, not a monotonicity
claim for the heat trace. If \(q=1\), every endomorphism is scalar,
and the heat trace is in fact independent of every metric.

\section{Finite telescoping and the direct multi-degree median bound}

For integers \(L>N\ge q-1\), all traces act on the same \(C\).
Applying VH12 at each step gives the finite total-variation bound
\begin{equation}\tag{VH13}\label{VH13}
\sum_{j=N}^{L-1}|\Psi_s(j+1)-\Psi_s(j)|
\le\frac1e\sum_{j=N}^{L-1}v_{j+1}
=\frac1e\log\frac{V_N}{V_L}.
\end{equation}
The endpoint difference is at most that sum. This is a finite statement;
it assumes no summability or decay along an infinite degree sequence.

Let \(U\) have columns \(i^{-j}b_j/\sqrt{\omega_j}\), \(N<j\le L\).
Every phase is retained in the column map. Repeated column addition
and direct multiplication give
\begin{equation}\tag{VH14}\label{VH14}
\begin{gathered}
G_L^{-1}=G_N^{-1}+UU^*,\quad
G_L=G_N-G_NU(I+U^*G_NU)^{-1}U^*G_N,\\
\frac{V_N}{V_L}=\det(I+U^*G_NU)
=\det(I+G_N^{1/2}UU^*G_N^{1/2}).
\end{gathered}
\end{equation}
For the determinant equality, eliminate either diagonal block of
\(\left(\begin{smallmatrix}I&U\\-U^*G_N&I\end{smallmatrix}\right)\).
This also proves the rectangular determinant identity without a rank
assumption. All matrices being inverted are positive.

Let \(\mu_1,\ldots,\mu_q\ge0\) be the eigenvalues, with all zeros
and multiplicities, of \(G_N^{1/2}UU^*G_N^{1/2}\), and put
\(l_j=\log(1+\mu_j)\). Then
\begin{equation}\tag{VH15}\label{VH15}
\boxed{\quad
|\Psi_s(L)-\Psi_s(N)|
\le\frac1e\min_{c\in\R}\sum_{j=1}^q|l_j-c|
\le\frac1e\sum_{j=1}^q l_j
=\frac1e\log\frac{V_N}{V_L}.\quad}
\end{equation}
Here any median of the complete list \(l_j\) realizes the minimum;
for even \(q\), the whole interval between its two middle entries
minimizes it. To prove the heat bound in full, put
\(C=G_N^{-1}G_L\), whose \(G_N\)-selfadjoint logarithm is
\(K=\log C\). Its eigenvalues are \(-l_j\), by VH14 and conjugation
with \(G_N^{1/2}\). Use the exact path
\(S_u=e^{uK/2}\), \(G_u=S_u^*G_NS_u\), \(M_u=S_uMS_u^{-1}\),
for \(0\le u\le1\). As in VH7 its physical map and unit are
\(\eta S_u^{-1}\) and \(S_u[1]\).
In the fixed \(G_N\)-adjoint, write \(H_u=M_u^{\dagger_N}M_u\).
Direct differentiation gives
\[
M_u'=\tfrac12[K,M_u],\quad
H_u'=M_u^{\dagger_N}KM_u-\tfrac12(KH_u+H_uK),
\]
and therefore
\[
\frac d{du}\tr e^{-sH_u}=-s\tr(KB_u),\quad
B_u=M_u e^{-sH_u}M_u^{\dagger_N}-H_u e^{-sH_u},\quad\tr B_u=0.
\]
The preceding singular-value argument gives \(\norm{B_u}_{G_N}\le1/(es)\).
For any real \(c\), the trace is unchanged on replacing \(K\) by
\(K-cI\). A singular-value expansion proves
\(|\tr XY|\le\norm X_{1,G_N}\norm Y_{G_N}\): sum its rank-one terms
and use Cauchy--Schwarz on each. Thus the derivative is bounded by
\(\norm{K-cI}_{1,G_N}/e\). Integrating proves VH15 after changing
the sign of \(c\). Finally \(\sum|l_j-c|\) has slope equal to the
number below \(c\) minus the number above it between consecutive
entries. This proves the median assertion, including repetitions.
If at most \(q/2\) entries are nonzero, zero is a median and the
volume bound agrees with this direct bound. In general the median
can improve it; if all entries are equal it gives zero exactly.

\section{The original Toda determinants and their complete phase term}

We now calculate the input of VH12 from the two original determinants,
not from a renamed volume sequence. Let
\begin{equation}\tag{VH16}\label{VH16}
\begin{gathered}
Z(\theta)=\int e^{\theta u}m_k(u)\,du=M_h(\theta)^k,\quad
Z_\chi(\theta)=\overline\chi(c-i\partial_\theta)
\chi(c+i\partial_\theta)Z(\theta)
=\int|\chi(c+iu)|^2d\mu_\theta(u),\\
\mathfrak D_n=\det[Z^{(i+j)}]_{i,j<n},\quad
\mathfrak B_n=\det[Z_\chi^{(i+j)}]_{i,j<n},\quad
\mathfrak D_0=\mathfrak B_0=1,\\
V_N=\frac{\mathfrak D_{N+1}}{\mathfrak B_{N-q+1}},\quad
\omega_j=\frac{\mathfrak D_{j+1}}{\mathfrak D_j},\quad
\nu_j=\frac{\mathfrak B_{j+1}}{\mathfrak B_j}.
\end{gathered}
\end{equation}
Differentiation under the integral proves the filter, with its two
signs and original \(c\). Its weight is the squared source norm of
the original relation, not its zero class after quotienting.
The basis \(1,S,\ldots,S^{q-1},\chi,\chi S,\ldots,\chi S^{N-q}\)
has triangular change of basis with diagonal one in the full source.
Its boundary Gram is the denominator in VH16; completing the square
in the boundary variables makes the Schur complement exactly \(G_N\).
The determinant of a block matrix with invertible boundary block is
the product of that determinant and its Schur complement, by triangular
elimination. This proves the volume ratio. The map from \(u^j\) to
\(S^j=(c+iu)^j\) has diagonal \(i^j\), whose determinant has modulus
one, so the two Gram determinants are unchanged by this exact phase
map. Triangular monic orthogonalization proves each norm ratio.

Use precisely the Toda conventions of \cite[(19), (25)--(26)]{TV}:
\begin{equation}\tag{VH17}\label{VH17}
\delta_j=V_j/V_{j-1},\quad v_j=-\log\delta_j\ (j\ge q),\quad
a_j=\omega_j/\omega_{j-1}\ (j\ge1),\quad
c_j=\nu_j/\nu_{j-1}\ (j\ge1),\quad c_0=0.
\end{equation}
Here \(a_j\) is the source recurrence coefficient, not the column
\(a\) in VH3. At every \(N\ge q-1\) the adjacent heat budget is
\begin{equation}\tag{VH18}\label{VH18}
v_{N+1}=\log\frac{\nu_{N-q+1}}{\omega_{N+1}},\qquad
|\Psi_s(N+1)-\Psi_s(N)|\le
\frac1e\log\frac{\nu_{N-q+1}}{\omega_{N+1}}.
\end{equation}
It is nonnegative because VH3 proves \(V_N/V_{N+1}\ge1\).
This compares the actual next boundary norm and the actual next
source norm, retaining their complete arithmetic mass.

The two Toda equations themselves are exact here:
\begin{equation}\tag{VH19}\label{VH19}
\begin{gathered}
(\log\mathfrak D_n)''=a_n,\quad (\log\mathfrak B_n)''=c_n,\quad
(\log V_N)''=a_{N+1}-c_{N-q+1},\\
v_{N+1}''=a_{N+1}-a_{N+2}-c_{N-q+1}+c_{N-q+2}.
\end{gathered}
\end{equation}
For a full proof, multiplication by \(u\) in the monic source
polynomials has recurrence
\(uU_j=U_{j+1}+b_j^{\rm rec}U_j+a_jU_{j-1}\), with \(a_0=0\).
All lower terms vanish by orthogonality and symmetry of multiplication;
pairing with \(U_{j-1}\) gives \(a_j=\omega_j/\omega_{j-1}\).
Differentiating orthogonality shows
\(\partial_\theta U_j=-a_jU_{j-1}\): the derivative has degree below
\(j\), and its pairing with every \(U_i\) except \(i=j-1\) vanishes,
while that pairing is \(-\omega_j\). Hence
\(\omega_j'/\omega_j=b_j^{\rm rec}\).
Differentiate \(\omega_j b_j^{\rm rec}=\int uU_j^2d\mu_\theta\),
insert the polynomial derivative and the recurrence, and obtain
\((b_j^{\rm rec})'=a_{j+1}-a_j\).
Since \(\mathfrak D_n=\prod_{j<n}\omega_j\), the sum telescopes to
the first equation. The boundary weight \(|\chi|^2d\mu_\theta\)
has the same exponential tilt and a fixed polynomial factor, so the
identical proof gives the second. Subtracting yields the last two
identities, with the zero-dimensional determinant handled by \(c_0=0\).
The last identity differentiates the exact budget function, not an
absolute-value heat inequality.

Here are the source-boundary contractions at the same degree:
\begin{equation}\tag{VH20}\label{VH20}
B_N=b_N^*G_Nb_N,\quad C_N=b_{N+1}^*G_Nb_{N+1},\quad
z_N=b_N^*G_Nb_{N+1},\quad
\ell_N=\log V_N,\quad \sigma=\tr M\in\R.
\end{equation}
Reflection of the original packet gives the real polynomial
\(i^{-q}\chi(c+iu)\); thus the companion \(M\) is conjugate to a real
matrix and its trace is real. Let \(\mathsf A_N\) be compression of
source multiplication by \(u\) through its least-norm quotient lifts.
It is \(G_N\)-selfadjoint. The outgoing column computation with the
retained phases gives
\begin{equation}\tag{VH21}\label{VH21}
\mathsf A_N=M+Q_N,\quad
Q_N=\frac{i}{\omega_N}b_{N+1}b_N^*G_N,\quad
\norm{Q_N}_{\mathrm{HS},G_N}^2=\frac{B_NC_N}{\omega_N^2},\quad
\frac{z_N}{\omega_N}=i(\sigma-\ell_N').
\end{equation}
For the first formula, full multiplication has only one degree beyond
\(N\); reduction of that column and the lifted last-coordinate row
has phase \(i^{-(N+1)}i^N=-i\), giving \(M=\mathsf A_N-Q_N\).
For the derivative identity, realize the minimizing lift as a polynomial
map \(\mathcal R_N:C\to\mathcal P_N\), with boundary space
\(\mathcal D_N=\chi\mathcal P_{N-q}\) and orthogonal projection
\(P_{\mathcal D_N}\) in \(L^2(\mu_\theta)\). Its class is fixed, so
\(\mathcal R_N'\) is a boundary. Differentiating its orthogonality
against a fixed boundary polynomial gives
\(\mathcal R_N'=-P_{\mathcal D_N}X\mathcal R_N\), where \(X\) is
multiplication by \(u\). Consequently
\(G_N'=\mathcal R_N^*X\mathcal R_N\): both lift-derivative terms vanish
against the orthogonal representatives. The compressed operator satisfies
\(G_N\mathsf A_N=\mathcal R_N^*X\mathcal R_N\), so
\(\ell_N'=\tr\mathsf A_N\), by differentiating the determinant.
Taking traces in the first formula of VH21 yields
\(\ell_N'=\sigma+i z_N/\omega_N\), which proves its last formula.
The squared norm follows by multiplying the rank-one operator and its
adjoint. In particular the phase in \(z_N\) is present at nonzero tilt.

For \(N\ge q\), the same preceding and next determinant updates give
\begin{equation}\tag{VH22}\label{VH22}
\frac{B_N}{\omega_N}=1-e^{-v_N},\quad
\frac{C_N}{\omega_{N+1}}=e^{v_{N+1}}-1,\quad
\norm{Q_N}_{\mathrm{HS},G_N}^2
=a_{N+1}(1-e^{-v_N})(e^{v_{N+1}}-1).
\end{equation}
Indeed the preceding update sends
\(b_N^*G_{N-1}b_N/\omega_N\) to that number divided by one plus itself,
by VH4. Its denominator is \(V_{N-1}/V_N=e^{v_N}\).
The next step is VH3. Multiplying gives the last identity.

The original rank-two allowance is not always this rank-one norm.
For the Hermitian defect \(W_N=i(Q_N^{\dagger_N}-Q_N)\), its two
possibly nonzero eigenvalues are \(\epsilon_N,-\epsilon_N\), with
\begin{equation}\tag{VH23}\label{VH23}
\boxed{\quad
\epsilon_N^2=\frac{B_NC_N-|z_N|^2}{\omega_N^2}
=a_{N+1}(1-e^{-v_N})(e^{v_{N+1}}-1)
-(\sigma-\ell_N')^2.\quad}
\end{equation}
Its trace is zero by VH21, since \(\tr Q_N=\ell_N'-\sigma\) is real.
It has rank at most two. The rank-one identity
\(Q_N^2=(\tr Q_N)Q_N\) and expansion give
\(\tr W_N^2=2\norm{Q_N}_{\mathrm{HS},G_N}^2-2(\tr Q_N)^2\).
A Hermitian rank-at-most-two trace-zero matrix has its two eigenvalues
opposite (both zero if its rank is below two), proving VH23.
For the original full reflected and conjugation-closed packet at
\(\theta=0\), parity makes \(z_N=0\), \(\sigma=\ell_N'(0)=0\),
and \(Q_N^2=0\). Only there do the two quantities coincide automatically.

When \(B_N>0\), an exact alternate form of the heat budget is
\begin{equation}\tag{VH24}\label{VH24}
v_{N+1}=\log\!\left(1+
\frac{\epsilon_N^2+(\sigma-\ell_N')^2}
{a_{N+1}(1-e^{-v_N})}\right),\qquad N\ge q.
\end{equation}
This is a substitution into the existing VH12 bound, not a new
asymptotic assertion. If \(B_N=0\), positivity gives \(b_N=0\),
\(Q_N=0\), and \(z_N=0\), but VH22 does not determine \(v_{N+1}\)
from those zeros. In that case use the undivided next contraction
\(C_N/\omega_{N+1}\) in VH3 instead; no zero denominator is introduced.
At the first admitted degree \(N=q-1\), the first \(q\) source columns
are a basis, so \(B_N/\omega_N=1\). The exact endpoint identities are
\begin{equation}\tag{VH25}\label{VH25}
v_q=\log(\nu_0/\omega_q),\quad
\norm{Q_{q-1}}_{\mathrm{HS},G_{q-1}}^2
=\frac{\nu_0-\omega_q}{\omega_{q-1}},\quad
\epsilon_{q-1}^2=\frac{\nu_0-\omega_q}{\omega_{q-1}}
-(\sigma-\ell_{q-1}')^2.
\end{equation}
These follow from VH18, VH21 and VH3 with leverage one; no inverse
at degree \(q-2\) is formed.

\section{A sharper scalar bound, with no rank multiplier}

Define \(d(1)=0\) and, for \(x>1\),
\(d(x)=(1-x^{-1})x^{-1/(x-1)}\). The rank-one nature of the actual
metric step permits the stronger result
\begin{equation}\tag{VH26}\label{VH26}
\boxed{\quad
|\Psi_s(N+1)-\Psi_s(N)|\le d(e^{v_{N+1}})
\le v_{N+1}/e.\quad}
\end{equation}
No factor \(q\) or \(\operatorname{rank}M\) occurs. The constant
\(d(\kappa)\) is sharp for arbitrary complex matrices under a prescribed
nonzero rank-one metric step in dimension at least two. This sharpness
statement does not claim that an arithmetic companion attains it.
The complete proof follows; an independent derivation and exact tests
are retained in \texttt{inequality\_check/RANK\_ONE\_HEAT\_BOUND.tex}.

First, for every measurable \(E\subset(0,\infty)\) with
\(\int_E dx/x\le\log\kappa\), \(\kappa>1\), one has
\begin{equation}\tag{VH27}\label{VH27}
\int_E e^{-x}\,dx\le d(\kappa).
\end{equation}
Put \(a=\log\kappa/(\kappa-1)\) and \(b=\kappa a\).
Integrating \(1/\kappa<1/t<1\) over \((1,\kappa)\) shows
\(0<a<1<b\). The function \(f(x)=xe^{-x}\) increases to \(1/e\)
at one and then decreases; moreover \(f(a)=f(b)\), because
\(b/a=\kappa\) and \(b-a=\log\kappa\).
Writing \(J=[a,b]\), \(c_*=f(a)\), and \(d\mu=dx/x\), gives
\[
\int_E f\,d\mu-\int_J f\,d\mu
=\int(1_E-1_J)(f-c_*)\,d\mu
+c_*[\mu(E)-\mu(J)]\le0.
\]
Both sets have finite \(\mu\)-measure, so the expression is integrable.
The first integrand is nonpositive both on \(J\) and its complement;
the second term is nonpositive. The remaining integral is
\(e^{-a}-e^{-b}=d(\kappa)\). This proves VH27 and also
\(d(\kappa)<\log\kappa/e\), since \(f<1/e\) except at one point.

Next let \(0\preceq X\preceq Y\preceq\kappa X\) be Hermitian
matrices with \(\operatorname{rank}(Y-X)\le1\). They have the same
kernel by the two quadratic-form inequalities. On its orthogonal
complement they are positive definite, and
\(Y-X=vv^*\). Thus
\(X^{-1/2}YX^{-1/2}=I+ww^*\preceq\kappa I\), with
\(w=X^{-1/2}v\), and
\(1\le\det Y/\det X=1+\norm w^2\le\kappa\) on that complement.
Order its positive eigenvalues as \(a_i\) for \(X\) and \(b_i\) for
\(Y\). The variational formula proved in the following sentence gives
\(a_i\le b_i\le a_{i+1}\), with the last upper inequality omitted.
For increasing eigenvalues the \(i\)-th is the minimum over
\(i\)-dimensional subspaces of the maximum Rayleigh quotient: the
first \(i\) eigenvectors attain it, while any such subspace intersects
the last \(r-i+1\) eigenvectors. The lower inequality follows from
\(Y\succeq X\); for the upper, intersect the first \(i+1\) eigenvectors
of \(X\) with \(v^\perp\), on which both quadratic forms agree.
The intervals \([sa_i,sb_i]\) have disjoint interiors and their union
has logarithmic measure \(\sum_i\log(b_i/a_i)\le\log\kappa\).
Their heat integral is the difference of the two traces. VH27 therefore
proves
\begin{equation}\tag{VH28}\label{VH28}
0\le\tr e^{-sX}-\tr e^{-sY}\le d(\kappa).
\end{equation}
The common kernel contributions cancel. Zero rank and \(\kappa=1\)
give equality zero directly and require no determinant on an empty
positive space.

Return to the actual \(G=G_N\), \(G_+=G_{N+1}\). Use the exact
isometries into Euclidean coordinates
\(S=G^{1/2}\) and \(DS\), where
\(h=Sa\), \(D=(I+hh^*)^{-1/2}\), and
\(G_+=S D^2 S\). Put \(A=SMS^{-1}\), \(B=DAD^{-1}\), \(C=DA\),
and \(F_s(Z)=\tr e^{-sZ^*Z}\). These maps carry the original physical
unit and observation by their inverses, exactly as in VH7, so
\(F_s(A)=\Psi_s(N)\), \(F_s(B)=\Psi_s(N+1)\).
For \(\kappa>1\), the orthogonal projection
\(P_h=hh^*/(\kappa-1)\) gives
\(D^2=I-(1-\kappa^{-1})P_h\), \(D^{-2}=I+(\kappa-1)P_h\).
Consequently
\begin{equation}\tag{VH29}\label{VH29}
\begin{gathered}
C^*C=A^*A-(1-\kappa^{-1})A^*P_hA,\quad
C^*C\preceq A^*A\preceq\kappa C^*C,\\
BB^*=CC^*+(\kappa-1)CP_hC^*,\quad
CC^*\preceq BB^*\preceq\kappa CC^*,\\
0\le F_s(C)-F_s(A)\le d(\kappa),\qquad
0\le F_s(C)-F_s(B)\le d(\kappa).
\end{gathered}
\end{equation}
Each difference is positive of rank at most one, so VH28 applies.
For the second application \(ZZ^*\) and \(Z^*Z\) have the same
positive eigenspaces via \(Z\), whose inverse on eigenvalue \(\lambda>0\)
is \(Z^*/\lambda\), and the same zero multiplicity \(q-\operatorname{rank}Z\).
Thus their heat traces agree including all zeros. Finally
\(F_s(B)-F_s(A)=[F_s(C)-F_s(A)]-[F_s(C)-F_s(B)]\) is the difference
of two numbers in \([0,d(\kappa)]\); its absolute value is at most
\(d(\kappa)\), not twice that number. This proves VH26 for arbitrary
complex \(M\); the zero update was already covered in VH6.

For sharpness with the actual prescribed metric step, let \(u=h/\norm h\)
and choose a Euclidean unit \(v\perp u\) in dimension at least two.
For any specified \(s>0\), choose
\begin{equation}\tag{VH30}\label{VH30}
M_+=S^{-1}(m_+uv^*)S,\quad
m_+^2=\frac{\kappa\log\kappa}{s(\kappa-1)},\qquad
M_-=S^{-1}(m_-vu^*)S,\quad
m_-^2=\frac{\log\kappa}{s(\kappa-1)}.
\end{equation}
The conjugation by \(D\) divides the nonzero squared singular value
of the first by \(\kappa\), and multiplies that of the second by
\(\kappa\). Their heat changes are respectively
\(+d(\kappa)\) and \(-d(\kappa)\), by direct substitution; their
\(q-1\) zero eigenvalues cancel. The complex direction of \(h\) is
kept in \(u\). These are sharp comparison operators, not substituted
arithmetic companions.

The original finite telescoping result can therefore be strengthened to
\begin{equation}\tag{VH31}\label{VH31}
|\Psi_s(L)-\Psi_s(N)|\le
\sum_{j=N}^{L-1}|\Psi_s(j+1)-\Psi_s(j)|
\le\sum_{j=N}^{L-1}d(e^{v_{j+1}})
\le\frac1e\log\frac{V_N}{V_L}.
\end{equation}
The minimum of its endpoint bound and the direct median bound VH15 is
also valid. VH18 and VH23--25 specify the exact original source and
boundary inputs to every step. No asymptotic estimate for those inputs,
no degree-uniform vanishing, and no sign of a particular arithmetic
heat change has been assumed or inferred.

\begin{thebibliography}{9}
\bibitem{TV} User-supplied programme continuation,
\emph{Source--boundary volume dynamics for the arithmetic control},
12 September 2026, \texttt{Tau\_Toda\_Volume\_Control/NOTE.tex},
equations (1)--(37), especially (16), (19), (23)--(29).
Complete original supplied TeX read for this calculation. This is
programme provenance, not a claim of human-paper authorship.
\bibitem{EM} Split-Zero programme,
\emph{Finite heat-trace control under an original quotient-metric comparison},
20 September 2026, \texttt{ENDPOINT\_METRIC\_HEAT.tex}, EM1--30.
Complete proof read; its log-metric bound is proved again here on the
actual degree metrics. Main045 HR1--32 and frozen043 NH1--36/HM1--19
remain unchanged.
\bibitem{Connes} Alain Connes, \emph{Heat expansion and zeta},
\url{https://arxiv.org/abs/2402.13082v1}, original author source
\texttt{heat-kernel-zeta.tex}, Gaussian formula \texttt{ft}.
Contextual citation inherited from the complete source reading recorded
by the preceding heat worker; not independently reread here. Its
RH-assuming infinite heat theorem is not used.
\bibitem{CV} Alain Connes and Walter D. van Suijlekom,
\emph{Quadratic Forms, Real Zeros and Echoes of the Spectral Action},
\url{https://arxiv.org/abs/2511.23257v1}, original source
\texttt{Araki-final-oct25.tex}, section \texttt{sect:sa}.
Context and reading boundary inherited from the preceding spectral-action
calculation. All finite identities used in this note are proved above.
\end{thebibliography}
\end{document}
